\chapter{Calculus and Optimisation}\label{ch:iv1:calculus-and-optimisation}

\begin{ivtrains}
Calculating with a clear head: integrals by symmetry, Lagrange multipliers, convexity, gradient and Newton steps, Stirling's formula, and the places where floating-point arithmetic gives wrong answers. The questions want a short route, and a check that the result is the right size.
\end{ivtrains}

\begin{question}[$\star$ \iqroles{QR, QT, Trader, Analyst} \iqfirm{any} \iqlevel{junior}]\label{iq:iv1:calculus-and-optimisation:1}
Where does $x^x$ reach its minimum for $x > 0$, and what is the minimum? Without a calculator, which is larger: $e^\pi$ or $\pi^e$? Prove it with a single function.
\end{question}

\begin{question}[$\star$ \iqroles{QR, QT, DS, QD} \iqfirm{any} \iqlevel{junior}]\label{iq:iv1:calculus-and-optimisation:2}
Compute $\int_{-\infty}^{\infty} e^{-x^2}\,dx$ by squaring it and using polar coordinates. Use the result to find the normal density's constant, and the mean of $|Z|$ for a standard normal $Z$.
\end{question}

\begin{question}[$\star$ \iqroles{QR, QT, QD, Dev} \iqfirm{any} \iqlevel{junior}]\label{iq:iv1:calculus-and-optimisation:3}
Use Newton's method on $f(x) = x^2 - 2$, starting from $x_0 = 1$. Write the update as a simple formula, and give the first three iterates as fractions. How many correct digits does each step add?
\end{question}

\begin{question}[$\star$ \iqroles{QR, QT, Trader, Analyst} \iqfirm{any} \iqlevel{junior}]\label{iq:iv1:calculus-and-optimisation:4}
Estimate $\ln 1.1$, $e^{0.1}$ and $\sqrt{1.04}$ to three decimals using only two or three terms of a Taylor series. A zero-coupon bond, paying 1 in 10 years, is priced at $e^{-10y}$, where $y$ is its yield. Using a first-order and a second-order approximation, what is the percentage change in price when the yield rises by one percentage point?
\end{question}

\begin{question}[$\star\star$ \iqroles{QR, QT, Trader} \iqfirm{any} \iqlevel{junior}]\label{iq:iv1:calculus-and-optimisation:5}
Compute $\int_0^{\pi/2} \frac{\sin^n x}{\sin^n x + \cos^n x}\,dx$ for any positive $n$. Then compute $\int_0^1 \frac{\ln(1 + x)}{1 + x^2}\,dx$, using a similar trick after the substitution $x = \tan\theta$.
\end{question}

\begin{question}[$\star\star$ \iqroles{QR, QT, DS, QD} \iqfirm{any} \iqlevel{junior}]\label{iq:iv1:calculus-and-optimisation:6}
Derive Stirling's formula for $n!$ from $n! = \int_0^\infty x^n e^{-x}\,dx$, by expanding around the peak of the integrand. Use it to approximate $\binom{100}{50}$. What is the chance of exactly 50 heads in 100 flips?
\end{question}

\begin{question}[$\star\star$ \iqroles{QR, QT, QD} \iqfirm{any} \iqlevel{junior}]\label{iq:iv1:calculus-and-optimisation:7}
Let $I_n = \int_0^{\pi/2} \sin^n x\,dx$. Derive a recursion using integration by parts, and compute $I_{10}$ exactly. How close is the large-$n$ approximation $\sqrt{\pi/(2n)}$?
\end{question}

\begin{question}[$\star\star$ \iqroles{QR, QT, Analyst, DS} \iqfirm{any} \iqlevel{junior}]\label{iq:iv1:calculus-and-optimisation:8}
Maximise $xyz$ subject to $x + y + z = 1$ with positive variables, by Lagrange multipliers and by another route. Then maximise $\sqrt{x} + \sqrt{y}$ subject to $x + y = B$. Show for $B = 2$ that the multiplier equals the change in the best value per unit of $B$, and say what this means.
\end{question}

\begin{question}[$\star\star$ \iqroles{QR, QT, DS, Analyst} \iqfirm{any} \iqlevel{junior}]\label{iq:iv1:calculus-and-optimisation:9}
Which of these are convex: $-\ln x$, $x^2$, $|x|$, $\sqrt{x}$, $\max(x, y)$, $xy$? Give a one-line reason for each. Use Jensen's inequality to explain why a normal annual return with mean 10\% and volatility 20\% compounds at a rate below 10\%, and estimate that rate.
\end{question}

\begin{question}[$\star\star$ \iqroles{QR, DS, QD} \iqfirm{any} \iqlevel{junior}]\label{iq:iv1:calculus-and-optimisation:10}
Gradient descent minimises $\frac12 x^\top Ax$ with $A$ having eigenvalues 1 and 100. What step sizes converge, and what is the best one? How many steps does it take to reduce the error by a factor of $10^6$? What does Newton's method do on this function, and what does heavy-ball momentum do?
\end{question}

\begin{question}[$\star\star$ \iqroles{QR, QD, Dev, QT} \iqfirm{any} \iqlevel{junior}]\label{iq:iv1:calculus-and-optimisation:11}
Give an example in which Newton's method on a root moves away from the root at each step, whatever the start. For $f(x) = \arctan x$, for what starting values does it fail? Why does solving for implied volatility need care, and how do you make a Newton iteration safe?
\end{question}

\begin{question}[$\star\star$ \iqroles{QR, QD, Dev, LLE} \iqfirm{any} \iqlevel{junior}]\label{iq:iv1:calculus-and-optimisation:12}
In double precision, the smaller root of $x^2 - 10^8x + 1 = 0$ from the usual formula comes out wrong. What does it give, what is the true value, and how do you fix it? Why does $1 - \cos x$ fail for small $x$, and what do you use instead? Why does the plain formula for variance fail on large offset values, and what is a better method?
\end{question}

\begin{question}[$\star\star$ \iqroles{QR, QD, Dev} \iqfirm{any} \iqlevel{junior}]\label{iq:iv1:calculus-and-optimisation:13}
You estimate a derivative with a forward difference in double precision. What step size is best, and what accuracy do you reach? How do both change for a central difference? Is there a way to reach almost full accuracy for an analytic function?
\end{question}

\begin{question}[$\star\star\star$ \iqroles{QR, QT} \iqfirm{hedge fund, market maker} \iqlevel{junior, experienced}]\label{iq:iv1:calculus-and-optimisation:14}
Compute $\int_0^{\pi/2} \ln \sin x\,dx$ with no series, using only symmetries of the sine and the double-angle formula.
\end{question}

\begin{question}[$\star\star\star$ \iqroles{QR, QT} \iqfirm{hedge fund} \iqlevel{junior, experienced}]\label{iq:iv1:calculus-and-optimisation:15}
Compute $\int_0^\infty \frac{\sin x}{x}\,dx$ by introducing a parameter: define $I(t) = \int_0^\infty e^{-tx}\frac{\sin x}{x}\,dx$, differentiate under the integral, and solve. What is $I(1)$? Why does the original integral converge, though not absolutely?
\end{question}

\begin{question}[$\star\star\star$ \iqroles{QR, QT, DS} \iqfirm{hedge fund} \iqlevel{junior, experienced}]\label{iq:iv1:calculus-and-optimisation:16}
Let $A$ be a symmetric positive definite $n$ by $n$ matrix. Compute $\int_{\mathbb{R}^n} e^{-\frac12 x^\top Ax}\,dx$. Deduce the covariance of the corresponding normal. Evaluate both for $A$ with rows $(2, 1)$ and $(1, 2)$. What is the sign of the correlation, and why?
\end{question}

\begin{question}[$\star\star\star$ \iqroles{QR, QT, DS, Risk} \iqfirm{hedge fund} \iqlevel{junior, experienced}]\label{iq:iv1:calculus-and-optimisation:17}
You allocate one unit of capital across three strategies, with amounts $x_i \ge 0$, to maximise $\sum_i \ln(1 + x_i/n_i)$, with $n = (0.2, 0.5, 1)$. Write the optimality conditions, with the sign constraints, and solve. Why is the last strategy unfunded? How does the answer change if the capital is 3 units?
\end{question}

\begin{question}[$\star\star\star$ \iqroles{QR, QD, DS} \iqfirm{any} \iqlevel{junior, experienced}]\label{iq:iv1:calculus-and-optimisation:18}
Prove that Newton's method converges quadratically near a simple root. Show that convergence is only linear at a double root, by using $f(x) = (x - 1)^2$, and find the rate. How would you fix the double-root case if you know the multiplicity?
\end{question}

\begin{question}[$\star\star\star\star$ \iqroles{QR, QT} \iqfirm{hedge fund} \iqlevel{experienced}]\label{iq:iv1:calculus-and-optimisation:19}
Show that $\sum_{k \ge 1} 1/k^2 = \pi^2/6$, by evaluating the double integral $\int_0^1\int_0^1 \frac{dx\,dy}{1 - xy}$ in two ways. Hint: expand as a geometric series for one route, and rotate the axes by 45 degrees for the other.
\end{question}

\begin{question}[$\star\star\star\star$ \iqroles{QR, QT} \iqfirm{hedge fund} \iqlevel{experienced}]\label{iq:iv1:calculus-and-optimisation:20}
Prove that $n! \sim C\sqrt{n}(n/e)^n$ for some constant $C$, by showing a suitable sequence is monotone and bounded. Then find $C = \sqrt{2\pi}$ exactly by using Wallis's result, $\lim_{n \to \infty} I_{2n}/I_{2n + 1} = 1$, with $I_n = \int_0^{\pi/2} \sin^n x\,dx$.
\end{question}
